\section{Time Complexity Analysis}
\label{subsec:complexity}

Let $V_{\mathrm{aff}}$ and $E_{\mathrm{aff}}$ denote vertices and edges in the structurally affected region,
$\Delta r = [\mathrm{rank}(\mathrm{SCC}(t)), \mathrm{rank}(\mathrm{SCC}(s))]$ the rank interval for Case~2,
$B$ the batch size, and $|\mathcal{B}_s|$ the number of reachable BiCC blocks.

\textbf{SPU (single-edge update).}
Table~\ref{tab:complexity_spu} gives the per-case cost.
Cases 1, 2 (no cycle), and 4 cost $\mathcal{O}(1)$.
All other cases are bounded by $\mathcal{O}(|V_{\mathrm{aff}}| + |E_{\mathrm{aff}}|)$,
confined to the affected structural region rather than the full graph.

\begin{table}[t]
\centering
\small
\setlength{\tabcolsep}{4pt}
\renewcommand{\arraystretch}{1.1}
\caption{SPU time complexity per case ($V_{\mathrm{aff}}, E_{\mathrm{aff}}$: affected region).}
\label{tab:complexity_spu}
\begin{tabular}{c l l}
\toprule
\textbf{Case} & \textbf{Operation} & \textbf{Complexity} \\
\midrule
1  & WCC merge                       & $\mathcal{O}(1)$ \\
2  & Inter-SCC, no cycle             & $\mathcal{O}(1)$ \\
2  & Inter-SCC, cycle                & $\mathcal{O}(\Delta r + |V_{\mathrm{aff}}| + |E_{\mathrm{aff}}|)$ \\
3  & Intra-SCC, inter-BiCC           & $\mathcal{O}(|V_{\mathrm{aff}}| + |E_{\mathrm{aff}}|)$ \\
4  & Intra-BiCC (no-op)              & $\mathcal{O}(1)$ \\
5a & Non-bridge deletion             & $\mathcal{O}(|V_{\mathrm{aff}}| + |E_{\mathrm{aff}}|)$ \\
5b & Bridge deletion within SCC      & $\mathcal{O}(|V_{\mathrm{aff}}| + |E_{\mathrm{aff}}|)$ \\
5c & Bridge deletion in DAG          & $\mathcal{O}(|V_{\mathrm{aff}}| + |E_{\mathrm{aff}}|)$ \\
\bottomrule
\end{tabular}
\end{table}

\textbf{FMU (batch update).}
Filter costs $\mathcal{O}(B)$; merge costs $\mathcal{O}(B \log B)$, producing $p \leq B$ independent regions.
The update stage runs SPU on each region concurrently, with total cost
$\mathcal{O}\!\left(B \log B + \max_{j}(|V_{\mathrm{aff},j}| + |E_{\mathrm{aff},j}|)\right)$.

\textbf{Two-level parallel execution.}
Level~1 traverses the condensation DAG in $\mathcal{O}(|\mathcal{V}_\mathcal{D}| + |\mathcal{E}_\mathcal{D}|)$
and the BC tree in $\mathcal{O}(|\mathcal{B}_s|)$.
Level~2 executes $|\mathcal{B}_s|$ tasks concurrently; parallel time is bounded by the slowest task.
Table~\ref{tab:complexity_parallel} compares GIR against sequential execution per algorithm.
The task-parallel speedup scales as 
$\frac{\sum_i T(S_i)}{\max_i T(S_i)}$.
reaching $|\mathcal{B}_s|$ when blocks are balanced.

\begin{table}[t]
\centering
\small
\setlength{\tabcolsep}{4pt}
\renewcommand{\arraystretch}{1.1}
\caption{Level~2 parallel vs.\ sequential complexity ($V_{B_i}, E_{B_i}$: vertices/edges in block $B_i$).}
\label{tab:complexity_parallel}
\begin{tabular}{l l l}
\toprule
\textbf{Algorithm} & \textbf{GIR (parallel)} & \textbf{Sequential} \\
\midrule
BFS/DFS & $\mathcal{O}\!\left(\max_i(|V_{B_i}|+|E_{B_i}|)\right)$
        & $\mathcal{O}(|V|+|E|)$ \\
SSSP    & $\mathcal{O}\!\left(\max_i(|E_{B_i}|\log|V_{B_i}|)\right)$
        & $\mathcal{O}(|E|\log|V|)$ \\
BC      & $\mathcal{O}\!\left(\max_i(|V_{B_i}||E_{B_i}|)\right)$
        & $\mathcal{O}(|V||E|)$ \\
\bottomrule
\end{tabular}
\end{table}
